\section{Proofs: groups of maps and of measure-space transformations}\label{sec:proofs-maps-transformations}

This section reproduces, from the preprint of the editor \cite{LupiniPointwise2026} in its version of 17 October 2026, the
proofs of the results stated in Section~\ref{sec:maps-transformations}. The text is that of the
preprint, with British spelling: its sections appear as subsections, its macros
are expanded and its labels renamed, and of its introduction only the
statements of the main results are kept. Citations of the companion preprints
point to their own proof pages where the cited statement is reproduced there.
The preprint has not been refereed. The last subsection derives the statements
of Section~\ref{sec:maps-transformations} from those of the preprint.

\subsection{General principles and pointwise completions}
\label{sec:pw-umf}

We recall the standard completion and extremely-amenable-kernel principles
used throughout; see \cite[Example~6.1.8 and Theorem~6.2.9]{Pestov2006} and
\cite[Lemma~2.3 and Corollary~2.5]{PestovRemarks2002}. Recall that \(G\) is \emph{extremely amenable} when \(M(G)\) is a
singleton.

\begin{lemma}\label{lem:pw-umf-principles}
Let \(G\) be a Hausdorff topological group.
\begin{enumerate}
\item If \(\iota:G\to Q\) is a continuous homomorphism with dense image in a
compact group \(Q\), then left translation through \(\iota\) makes \(Q\) a
minimal \(G\)-flow. Consequently, there is a factor map
\[
M(G)\longrightarrow Q.
\]
\item If \(\iota\) is a topological embedding, so that \(Q\) is the compact
group completion of the precompact group \(G\), then
\[
M(G)\cong Q.
\]
\item If \(N\) is a closed extremely amenable normal subgroup of \(G\), then
\[
M(G)\cong M(G/N).
\]
In particular, if \(G/N\) is compact, then \(M(G)\cong G/N\).
\end{enumerate}
\end{lemma}

\begin{proof}
For the first assertion, every \(G\)-orbit in \(Q\) is dense because
\(\iota(G)\) is dense. The universal property of \(M(G)\) gives the factor
map.

For the second, every compact \(G\)-flow extends uniquely to the completion
\(Q\); see \cite[Fact~2.3]{BassoZucker2021}. If \(X\) is minimal and \(x\in X\), the
extended orbit map \(Q\to X\), \(q\mapsto qx\), is onto. Thus the minimal
\(G\)-flow \(Q\) maps onto every minimal \(G\)-flow.

Finally, let \(X\) be a minimal \(G\)-flow. Extreme amenability supplies an
\(N\)-fixed point. Since \(N\) is normal, the fixed-point set \(X^N\) is
nonempty, closed, and \(G\)-invariant, and hence \(X^N=X\). Thus every
minimal \(G\)-flow factors through \(G/N\), which proves the last assertion.
\end{proof}


\subsubsection{Finite interpolation beyond manifolds}

The pointwise-completion argument only needs the ability to realise arbitrary
finite prescriptions.  This gives a general function-space statement.  We use
the standard notation \(C_p(X,K)\) for the group of continuous maps endowed
with pointwise convergence. The separation property in
\cite[Definition~2.1 and Remark~2.2]{HernandezRodenas2007} supplies finite interpolation;
\cite[Definition~2.2(iii), Proposition~2.3(i), and proof of Theorem~3.3,
equation~(10)]{ShakhmatovSpevak2010} records the resulting completion. The next statement is
its universal-minimal-flow consequence, including the relative version.

\begin{theorem}\label{thm:pw-cp-umf}
Let \(X\) be a Tychonoff space and let \(K\) be a compact path-connected
topological group. Then
\[
M(C_p(X,K))\cong K^X.
\]
More generally, if \(A\subseteq X\) is closed and
\[
C_p(X,A;K)=\{f\in C(X,K):f|_A=e\},
\]
with its pointwise topology, then
\[
M(C_p(X,A;K))\cong K^{X\setminus A}.
\]
In the second formula the omitted coordinates on \(A\) are fixed at the
identity.
\end{theorem}

\begin{proof}
\uses{lem:pw-umf-principles}
Fix distinct points \(x_1,\ldots,x_m\in X\) and prescribed values
\(k_1,\ldots,k_m\in K\). Complete regularity supplies continuous functions
\(\psi_i:X\to[0,1]\) such that
\(\psi_i(x_i)=1\) and \(\psi_i(x_j)=0\) for \(j\ne i\). Choose paths
\(c_i:[0,1]\to K\) from \(e\) to \(k_i\). Then
\[
f(x)=\prod_{i=1}^m c_i(\psi_i(x))
\]
is continuous and satisfies \(f(x_i)=k_i\). Thus the evaluation image of
\(C_p(X,K)\) is dense in \(K^X\).

For the relative group, take the separating functions \(\psi_i\) to vanish on
the closed set \(A\) as well as at the other prescribed points. The same
construction gives density in the closed subgroup of \(K^X\) whose coordinates
on \(A\) are the identity. In both cases this is precisely the compact group
completion, so Lemma~\ref{lem:pw-umf-principles} applies.
\end{proof}

\medskip
The connectedness assumption can be removed if one records the Boolean
information carried by the components. For a Tychonoff space \(X\), write
\(x\sim_{\mathrm{clop}}y\) when every clopen subset of \(X\) contains either
both \(x,y\) or neither of them. Thus the equivalence classes are the
quasi-components of \(X\). If \(A\subseteq X\), put
\[
 \operatorname{bcl}_X(A)
 =\bigcap\{U\subseteq X:U\text{ is clopen and }A\subseteq U\},
\]
with \(\operatorname{bcl}_X(\varnothing)=\varnothing\). We call this the
Boolean hull of \(A\).

Let \(K\) be a compact Lie group, let \(q_0:K\to K/K_0\) be the component
quotient, and define
\begin{align*}
 \mathcal B_X(K)
 &=\{u\in K^X:q_0(u(x))=q_0(u(y))
          \text{ whenever }x\sim_{\mathrm{clop}}y\},\\
 \mathcal B_{X,A}(K)
 &=\{u\in\mathcal B_X(K):u|_A=e,\quad
       q_0(u(x))=K_0\text{ for }x\in\operatorname{bcl}_X(A)\}.
\end{align*}

\begin{theorem}\label{thm:pw-cp-disconnected-umf}
Let \(X\) be Tychonoff and let \(K\) be a compact Lie group, not necessarily
connected. Then
\[
 M(C_p(X,K))\cong\mathcal B_X(K).
\]
If \(A\subseteq X\) is closed, then
\[
 M(C_p(X,A;K))\cong\mathcal B_{X,A}(K).
\]
In particular, the completion is \(K^X\) whenever the quasi-components of
\(X\) are singletons, and it is the subgroup of \(K^X\) with globally
constant component coordinate whenever \(X\) is connected.
\end{theorem}

\begin{proof}
\uses{lem:pw-umf-principles}
Every continuous map from \(X\) to the finite discrete group \(K/K_0\) is
locally constant. Its pointwise closure consists exactly of the functions
\(c:X\to K/K_0\) which are constant on the
\(\sim_{\mathrm{clop}}\)-classes.
Indeed, finitely many distinct classes can be separated by a finite clopen
partition. In the relative case the locally constant maps which equal \(K_0\)
on \(A\) have as their pointwise closure the functions which, in addition,
equal \(K_0\) on \(\operatorname{bcl}_X(A)\). To see the converse, separate
each of the finitely many points outside the Boolean hull from \(A\) by a
clopen set, and refine these sets together with the pairwise separating
clopen sets to a finite clopen partition.

It remains to interpolate inside the components. Fix finitely many points
\(x_1,\ldots,x_m\) and target values prescribed by an element of the displayed
compact group. Choose a locally constant component map with the required
values and, in the relative case, equal to \(K_0\) on \(A\). Choose one
representative of every component of \(K\), taking the representative of
\(K_0\) to be \(e\); composition gives a continuous preliminary map
\(f_0:X\to K\). For each \(x_i\notin A\), the correction
\(f_0(x_i)^{-1}u(x_i)\) lies in the connected compact Lie group \(K_0\).
Choose a path in \(K_0\) from \(e\) to this correction and a continuous scalar
function which is one at \(x_i\) and zero on \(A\) and at the other prescribed
points. Multiplying the resulting corrections into \(f_0\) gives a continuous
map, equal to \(e\) on \(A\) in the relative case, with all the prescribed
values. Thus the evaluation images are dense in the stated compact groups,
which are their compact group completions. Apply
Lemma~\ref{lem:pw-umf-principles}.
\end{proof}

There is also a projective version. If \(Z\leq K\) is closed and central,
write \(\Delta Z\) for the compact subgroup of constant \(Z\)-valued maps.

\begin{corollary}\label{cor:pw-projective-cp-umf}
With the quotient of the pointwise topology,
\[
 M\bigl(C_p(X,K)/\Delta Z\bigr)
 \cong \mathcal B_X(K)/\Delta Z.
\]
For connected \(K\), this reduces to \(K^X/\Delta Z\).
\end{corollary}

\begin{proof}
\uses{lem:pw-umf-principles, thm:pw-cp-disconnected-umf}
The completion of the quotient by the compact subgroup \(\Delta Z\) is the
quotient of the compact completion by \(\Delta Z\). The result follows from
Theorem~\ref{thm:pw-cp-disconnected-umf} and
Lemma~\ref{lem:pw-umf-principles}.
\end{proof}

There is a bundle version. Let \(\mathcal H\to X\) be a locally trivial bundle
of topological groups with typical fibre \(K\), and let
\(\Gamma_p(X,\mathcal H)\) be its continuous-section group with pointwise
convergence.

\begin{corollary}\label{cor:pw-topological-section-umf}
Let \(X\) be paracompact Hausdorff and let \(\mathcal H\to X\) be a locally
trivial bundle of compact path-connected groups. Then
\[
M(\Gamma_p(X,\mathcal H))
 \cong\prod_{x\in X}\mathcal H_x.
\]
If \(A\subseteq X\) is closed, the section group fixed at the identity on
\(A\) has universal minimal flow
\[
\prod_{x\in X\setminus A}\mathcal H_x.
\]
\end{corollary}

\begin{proof}
\uses{lem:pw-umf-principles}
Given finitely many base points, choose pairwise disjoint trivializing
neighbourhoods. In each neighbourhood, a scalar function which is one at the
chosen point and zero off the neighbourhood, composed with a path in the fibre
from the identity to the prescribed value, gives a local section which patches
with the identity section. Multiplying the finitely many local sections proves
density in the displayed fibre product. If the section must be the identity on
\(A\), choose the scalar functions to vanish on \(A\). Apply
Lemma~\ref{lem:pw-umf-principles}.
\end{proof}

The finite-interpolation principle also applies to a regularity class which
does not contain the smooth maps.

\begin{corollary}\label{cor:pw-analytic-mapping-umf}
Let \(X\) be a compact real-analytic manifold and let \(K\) be a connected
compact Lie group. With the topology of pointwise convergence,
\[
M(C^\omega(X,K))\cong K^X.
\]
\end{corollary}

\begin{proof}
\uses{lem:pw-umf-principles}
By the real-analytic embedding theorem \cite{Grauert1958}, choose a real-analytic
embedding \(j:X\to\mathbb R^N\). For distinct
\(x_1,\ldots,x_m\), ordinary polynomial interpolation on the finite subset
\(\{j(x_1),\ldots,j(x_m)\}\) gives polynomials \(p_i\) with
\(p_i(j(x_r))=\delta_{ir}\). Given \(k_i\in K\), choose
\(A_i\in\mathfrak k\) with \(\exp(A_i)=k_i\); the exponential map is
surjective for connected compact Lie groups. The real-analytic map
\[
x\longmapsto\prod_{i=1}^m
 \exp\bigl(p_i(j(x))A_i\bigr)
\]
takes the prescribed values. Hence \(C^\omega(X,K)\) is dense in \(K^X\),
and the conclusion follows from Lemma~\ref{lem:pw-umf-principles}.
\end{proof}

The same argument applies to metric regularity classes which permit finite
scalar interpolation. Fix a compatible bi-invariant Riemannian metric on a
compact connected Lie group \(K\). For a metric space \(X\), write
\(\operatorname{Lip}_p(X,K)\) for all Lipschitz maps with pointwise
convergence, and, for \(0<\gamma\leq1\), define
\(C_p^{0,\gamma}(X,K)\) analogously.

\begin{corollary}\label{cor:pw-lipschitz-mapping-umf}
Let \(X\) be any metric space and let \(K\) be a compact connected Lie group.
Then
\[
 M(\operatorname{Lip}_p(X,K))\cong K^X,
 \qquad
 M(C_p^{0,\gamma}(X,K))\cong K^X
 \quad(0<\gamma\leq1).
\]
\end{corollary}

\begin{proof}
\uses{lem:pw-umf-principles}
On a finite subset \(\{x_1,\ldots,x_m\}\subseteq X\), choose bounded
Lipschitz functions \(\psi_i:X\to[0,1]\) with
\(\psi_i(x_j)=\delta_{ij}\). Given \(k_i\in K\), choose
\(A_i\in\mathfrak k\) with \(\exp(A_i)=k_i\). Then
\[
 x\longmapsto\prod_{i=1}^m\exp(\psi_i(x)A_i)
\]
is Lipschitz and takes the prescribed values. A bounded Lipschitz map is
\(\gamma\)-H\"older for every \(0<\gamma\leq1\), after changing the constant.
Thus both evaluation images are dense in \(K^X\), and
Lemma~\ref{lem:pw-umf-principles} applies.
\end{proof}

There is also a topological compact-support version, extending the smooth
manifold statement above.  Write \(C_{c,p}(X,K)\) for the group of continuous
maps \(f:X\to K\) that equal the identity outside a compact subset of \(X\),
equipped with the topology of pointwise convergence.

\begin{corollary}\label{cor:pw-lch-compact-support-umf}
Let \(X\) be locally compact Hausdorff and let \(K\) be a compact
path-connected topological group. With pointwise convergence,
\[
 M(C_{c,p}(X,K))\cong K^X.
\]
\end{corollary}

\begin{proof}
\uses{thm:pw-cp-umf}
Given finitely many points, local compactness and complete regularity provide
scalar separating functions with compact support. Composing them with paths
in \(K\) and multiplying as in Theorem~\ref{thm:pw-cp-umf} realises every finite
prescription. Hence the compact completion is \(K^X\).
\end{proof}

Finite interpolation has a holomorphic form as well. If \(S\) is a Stein
manifold and \(T\) is a complex Lie group, write
\(\mathcal O_p(S,T)\) for the holomorphic mapping group with pointwise
convergence.

\begin{theorem}\label{thm:pw-holomorphic-mapping-umf}
Let \(S\) be a Stein manifold and let \(T\) be a connected compact complex Lie
group. Then
\[
 M(\mathcal O_p(S,T))\cong T^S.
\]
If \(Z\leq T\) is closed and central, then
\[
 M(\mathcal O_p(S,T)/\Delta Z)\cong T^S/\Delta Z.
\]
\end{theorem}

\begin{proof}
\uses{lem:pw-umf-principles}
A connected compact complex Lie group \(T\) is abelian. Indeed, the
holomorphic matrix entries of its adjoint representation are constant by
compactness, so its Lie algebra is abelian; connectedness then gives the
assertion for \(T\). Its exponential map is onto, since its image is an open
subgroup. Fix distinct \(x_1,\ldots,x_m\in S\). Holomorphic separability of
the Stein manifold supplies functions \(h_j\) with
\(h_j(x_i)=\delta_{ij}\): for each \(i\ne j\), normalise a separating
function to be zero at \(x_i\) and one at \(x_j\), and take the product over
\(i\). Given \(t_j\in T\), choose \(v_j\) in its Lie algebra with
\(\exp(v_j)=t_j\). Then
\[
 x\longmapsto\exp\!\left(\sum_{j=1}^m h_j(x)v_j\right)
\]
is a holomorphic interpolant. Thus \(\mathcal O_p(S,T)\) is dense in
\(T^S\), and Lemma~\ref{lem:pw-umf-principles} and the compact central
quotient give the two formulas. The Oka interpolation theorem
\cite[Theorem~1.1]{ForstnericPrezelj2001} provides a more general route, but is unnecessary
under the present compact-target hypotheses.
\end{proof}

These regularity examples are applications of finite interpolation and the
standard precompact-flow calculation. In particular, the continuous case
uses the explicit earlier completion in \cite{ShakhmatovSpevak2010}; the other displayed
completions are derived directly, without a claim of a new interpolation
principle or historical priority for every formulation.


\subsubsection{One-dimensional NCCW unitary groups}

Let \(E\) and \(F\) be finite-dimensional unital \(C^*\)-algebras, let
\[
\phi=(\phi_0,\phi_1):E\longrightarrow F\oplus F
\]
be a unital \(*\)-homomorphism, and consider the one-dimensional NCCW algebra
\[
A=\{(p,u)\in C(I,F)\oplus E:(p(0),p(1))=\phi(u)\}.
\]
Write
\[
\phi_{\mathrm U}=\phi|_{\mathrm U(E)}:
\mathrm U(E)\longrightarrow\mathrm U(F)\times\mathrm U(F),
\]
and put
\begin{equation}\label{eq:pw-nccw-pointwise-completion}
\mathcal B_A=
\left\{(x,u)\in\mathrm U(F)^I\times\mathrm U(E):
 (x(0),x(1))=\phi_{\mathrm U}(u)\right\}.
\end{equation}
Equip \(\mathrm U(A)\) with the pointwise topology induced by the
\(\mathrm U(E)\)-coordinate and by all evaluations of the path coordinate.

\begin{proposition}\label{prop:pw-nccw-umf}
The compact group \(\mathcal B_A\) is the group completion and universal
minimal flow of pointwise-topological \(\mathrm U(A)\):
\[
M(\mathrm U(A),\tau_{\mathrm{pw}})\cong\mathcal B_A.
\]
For the norm topology, \(\mathcal B_A\) is a compact minimal
\(\mathrm U(A)\)-flow, and hence there is a factor map
\[
M(\mathrm U(A)_{\|\cdot\|})\longrightarrow\mathcal B_A.
\]
If the fibre algebra \(F\) is nonzero, as in the standard unital
one-dimensional NCCW convention, then
\(M(\mathrm U(A)_{\|\cdot\|})\) is nonmetrizable.
\end{proposition}

\begin{proof}
\uses{lem:pw-umf-principles}
The space \(\mathcal B_A\) is the fibre product of the endpoint homomorphism
\(\mathrm U(F)^I\to\mathrm U(F)\times\mathrm U(F)\) and
\(\phi_{\mathrm U}\), so it is a closed subgroup of the compact group
\(\mathrm U(F)^I\times\mathrm U(E)\).

The finite-dimensional group \(\mathrm U(F)\) is a finite product of connected
unitary groups. Given \((x,u)\in\mathcal B_A\) and finitely many coordinates,
add \(0\) and \(1\) to the coordinate set, and join the prescribed values in
their order by continuous paths in \(\mathrm U(F)\), retaining the endpoint
values prescribed by \(\phi_{\mathrm U}(u)\). This gives
\((p,u)\in\mathrm U(A)\) agreeing with \((x,u)\) at the selected coordinates.
Hence \(\mathrm U(A)\) is dense in \(\mathcal B_A\), and
Lemma~\ref{lem:pw-umf-principles} proves both the pointwise calculation and the
norm-topological factor assertion.

Taking \(u=1_E\) and fixing the endpoint coordinates at \(1_F\) gives a closed
copy of \(\mathrm U(F)^{(0,1)}\) in \(\mathcal B_A\). Since \(F\) is a
nonzero unital finite-dimensional \(C^*\)-algebra under the stated
convention, \(\mathrm U(F)\) is
nontrivial, and this product is nonmetrizable. Since a Hausdorff continuous
image of a compact metrizable space is metrizable, the displayed factor map
forces \(M(\mathrm U(A)_{\|\cdot\|})\) to be nonmetrizable.
\end{proof}

\begin{remark}\label{rem:pw-nccw-pointwise-coordinate}
\uses{prop:pw-nccw-umf}
Including the \(\mathrm U(E)\)-coordinate in \(\tau_{\mathrm{pw}}\) is
essential. If the topology is induced only by evaluations of the path
coordinate, then the compact kernel \(\ker(\phi_{\mathrm U})\) is not
separated; the Hausdorff quotient instead completes to
\[
\{x\in\mathrm U(F)^I:
  (x(0),x(1))\in\phi_{\mathrm U}(\mathrm U(E))\}.
\]
Proposition~\ref{prop:pw-nccw-umf} identifies \(\mathcal B_A\) only as a
factor in the norm topology, not as the full norm-topological universal
minimal flow.
\end{remark}


\subsection{Convergence in measure and sampled topologies}

The opposite extreme from pointwise convergence is provided by measurable
maps. Let \((Y,\mu)\) be a standard probability space and let
\[
A_\mu=\{a\in Y:\mu(\{a\})>0\}
\]
be its at most countable set of atomic points.

\begin{proposition}\label{prop:pw-l0-umf}
For every compact metrizable group \(K\), with \(L^0(Y,\mu;K)\) carrying
convergence in measure,
\[
M(L^0(Y,\mu;K))\cong K^{A_\mu}.
\]
In particular, \(L^0(Y,\mu;K)\) is extremely amenable when \(\mu\) is
atomless.
\end{proposition}

\begin{proof}
\uses{lem:pw-umf-principles}
Let \(Y_{\mathrm c}=Y\setminus A_\mu\). Restriction gives a topological group
isomorphism
\[
L^0(Y,\mu;K)
\cong L^0(Y_{\mathrm c},\mu;K)\times K^{A_\mu}.
\]
Indeed, for a bounded compatible bi-invariant metric \(d_K\), the atomic
factor carries the weighted metric
\[
\sum_{a\in A_\mu}\mu(\{a\})d_K(k_a,\ell_a),
\]
which induces the product topology. For the compact target used here, extreme amenability of the atomless
\(L^0\)-group is already \cite[Theorem~2.20]{GiordanoPestov2007}; the circle case appears
in \cite[Theorems~1.2--1.3]{Glasner1998}. Pestov and Schneider prove the more
general amenable-target theorem \cite[Theorem~1.1]{PestovSchneider2017}. If \(\mu(Y_{\mathrm c})>0\), normalise the
restriction of \(\mu\) to \(Y_{\mathrm c}\), which does not change the
convergence-in-measure topology; the resulting standard atomless probability
space is isomorphic modulo null sets to the standard model. If
\(\mu(Y_{\mathrm c})=0\), the first factor is trivial. The compact group \(K\)
is amenable, so Lemma~\ref{lem:pw-umf-principles} applies to the projection onto
\(K^{A_\mu}\).
\end{proof}

The same argument permits the compact group to vary measurably.  Say that a
Borel bundle of compact metrizable groups \(\mathcal H\to Y\) is
\emph{countably trivializable} if there is a countable Borel partition
\(Y=\bigsqcup_jY_j\) and compact metrizable groups \(H_j\) such that
\(\mathcal H|_{Y_j}\) is Borel-isomorphic, as a group bundle, to
\(Y_j\times H_j\).

\begin{theorem}\label{thm:pw-measurable-bundle-umf}
Let \((Y,\mu)\) be a standard probability space and let
\(\mathcal H\to Y\) be a countably trivializable Borel bundle of compact
metrizable groups.  Equip the measurable-section group
\(L^0(Y,\mu;\mathcal H)\) with convergence in measure. Then
\[
M(L^0(Y,\mu;\mathcal H))
 \cong\prod_{a\in A_\mu}\mathcal H_a.
\]
In particular, this measurable-section group is extremely amenable when
\(\mu\) is atomless.
\end{theorem}

\begin{proof}
\uses{lem:pw-umf-principles}
Restriction to the atoms is a continuous surjective homomorphism onto the
displayed compact product. Its kernel is the measurable-section group over
\(Y\setminus A_\mu\). After applying the countable trivializing partition,
this kernel is a product of groups of the form
\(L^0(Y_j\setminus A_\mu,\mu;H_j)\). Each nontrivial factor is extremely
amenable by \cite[Theorem~1.1]{PestovSchneider2017}, and products of extremely amenable
groups are extremely amenable. Here convergence in measure gives the product
topology on the partitioned group: after finitely many pieces the remaining
total measure is arbitrarily small. Lemma~\ref{lem:pw-umf-principles} now
identifies the universal minimal flow with the compact quotient.
\end{proof}

\begin{corollary}\label{cor:pw-measurable-gauge-umf}
Let \(Y\) be a compact manifold with a regular Borel probability measure
\(\mu\), let \(K\) be a compact Lie group, and let \(P\to Y\) be a smooth
principal \(K\)-bundle.  Put
\[
\operatorname{Gau}^0_\mu(P)=L^0(Y,\mu;\operatorname{Ad}P)
\]
be its measurable gauge group, modulo equality almost everywhere, with
convergence in measure. Then
\[
M(\operatorname{Gau}^0_\mu(P))
\cong\prod_{a\in A_\mu}\operatorname{Ad}(P)_a.
\]
In particular, the measurable gauge group is extremely amenable when
\(\mu\) is atomless.
\end{corollary}

\begin{proof}
\uses{thm:pw-measurable-bundle-umf}
A finite or countable Borel partition subordinate to bundle trivializations
shows that \(\operatorname{Ad}(P)\) is countably trivializable. Apply
Theorem~\ref{thm:pw-measurable-bundle-umf}.
\end{proof}

\begin{corollary}\label{cor:pw-measure-topology-umf}
Suppose that \(K\) is connected. Equip any path, loop, twisted loop, finite
endpoint-pullback, or one-dimensional gauge group considered above with the
coarser topology induced solely by convergence of its map or section
coordinate in normalised volume measure. If this topology does not separate
auxiliary coordinates of a pullback, take its Hausdorff quotient. Then the
resulting group is extremely amenable, and its universal minimal flow is a
point. The same conclusion holds after imposing finitely many basepoint or
framing conditions.
\end{corollary}

\begin{proof}
\uses{prop:pw-l0-umf}
The map into the appropriate measurable-map or measurable-section group
induces precisely the stated topology on its Hausdorff image. That image is
dense. For maps, approximate a measurable map by a finite-valued simple map,
approximate its level sets by finite unions of intervals, and join the
finitely many values across transition intervals of arbitrarily small total
measure. For gauge groups, cut the one-manifold into finitely many trivializing
arcs and place the transitions in neighbourhoods of arbitrarily small total
measure; equivalently, use a finite Borel partition subordinate to bundle
trivializations to identify the measurable-section group with an
\(L^0\)-group. Connectedness of \(K\) supplies the joining paths. Endpoint,
derivative-matching, twisting, and framing conditions can be repaired on
neighbourhoods whose measures tend to zero. Thus the completion is an
atomless \(L^0\)-group, which is extremely amenable by
Proposition~\ref{prop:pw-l0-umf}. Universal minimal flows are unchanged on
passing to a dense completion by \cite[Fact~2.3]{BassoZucker2021}.
\end{proof}

Only the map or section coordinate enters this corollary. For example, in an
NCCW pullback a unitary vertex coordinate need not be determined by the path
almost everywhere. If that coordinate is retained with its compact topology,
the completion instead retains it: for connected unitary fibres it is
\(L^0(I;\mathrm U(F))\times\mathrm U(E)\), and its universal minimal flow
is \(\mathrm U(E)\). Endpoint conditions can be repaired on intervals of
arbitrarily small measure, so they impose no relation between these two
completion coordinates. This is a different topology from the one in the
corollary.

\begin{remark}\label{rem:pw-disconnected-measure-umf}
\uses{cor:pw-measure-topology-umf}
For disconnected \(K\), the component restriction does not disappear under
completion. We use the same Hausdorff-quotient convention for auxiliary
coordinates as in Corollary~\ref{cor:pw-measure-topology-umf}. Let \(C=K/K_0\), let \(\overline L\) be the image of
\(L\leq K\times K\) in \(C\times C\), and set
\[
C_L=\{c\in C:(c,c)\in\overline L\}.
\]
For a finite endpoint pullback with convergence in Lebesgue measure, the
completion has an extremely amenable normal \(L^0(I,K_0)\)-kernel and compact
quotient \(C_L\), and hence its universal minimal flow is \(C_L\). By
\cite[Fact~2.3]{BassoZucker2021}, the original dense endpoint pullback has the same
universal minimal flow. Thus unbased paths and ordinary loops give \(C\),
based groups give a point, and an \(\alpha\)-twisted loop group gives
\(\operatorname{Fix}(\pi_0(\alpha))\).
\end{remark}


\subsubsection{Mixed sampling topologies and measure-space symmetries}

There is a continuum of useful group topologies between convergence in measure
and evaluation at every point. Let \(X\) be a compact connected smooth
manifold, let \(\mu\) be an atomless regular Borel probability measure of full
support, let \(K\) be a connected compact Lie group, and fix an arbitrary set
\(D\subseteq X\). Define \(\tau_{\mu,D}\) to be the topology on
\(C^\infty(X,K)\) induced by
\[
\jmath_D:C^\infty(X,K)\longrightarrow
 L^0(X,\mu;K)\times K^D,
 \qquad f\longmapsto([f],(f(d))_{d\in D}).
\]
Full support makes this homomorphism injective.

\begin{theorem}\label{thm:pw-mixed-sampling-umf}
For every \(D\subseteq X\),
\[
M(C^\infty(X,K),\tau_{\mu,D})\cong K^D.
\]
Thus \(D=\varnothing\) gives extreme amenability, while for \(D=X\) the
universal minimal flow agrees with the one for the pointwise topology.
\end{theorem}

\begin{proof}
\uses{lem:pw-umf-principles}
Smooth maps are dense in \(L^0(X,\mu;K)\). One way to see this is to
approximate a measurable map by a finite-valued simple map, approximate its
level sets inside finitely many coordinate charts, and insert smooth transition
regions of arbitrarily small measure; connectedness of \(K\) supplies paths
between the finitely many values.

This approximation can be made simultaneously with any finite prescription on
\(D\). After first choosing a smooth approximation, multiply it by corrections
supported in pairwise disjoint coordinate balls around the finitely many sampled
points. The balls can be chosen with arbitrarily small total measure, so the
corrections impose the prescribed point values without changing the limit in
measure. Hence \(\jmath_D\) has dense image and the two-sided completion is
\[
L^0(X,\mu;K)\times K^D.
\]
The first factor is extremely amenable by \cite[Theorem~1.1]{PestovSchneider2017}; it is a
closed normal subgroup and the quotient is the compact group \(K^D\).
Lemma~\ref{lem:pw-umf-principles} gives the formula.
\end{proof}

The Hausdorff completion underlying the preceding result can be computed for
arbitrary measurable maps, without connectedness or smoothness assumptions.
Let \(\mathcal M(Y,K)\) be the group of Borel maps \(Y\to K\), with no
almost-everywhere identification. For \(D\subseteq Y\), let
\[
 \jmath_{\mu,D}:\mathcal M(Y,K)\longrightarrow
 L^0(Y,\mu;K)\times K^D,
 \qquad f\longmapsto([f],f|_D),
\]
and write \(G_{\mu,D}(K)\) for its image with the subspace topology.
Equivalently, this is the Hausdorff quotient of \(\mathcal M(Y,K)\) by maps
which equal the identity almost everywhere and at every point of \(D\).

\begin{theorem}\label{thm:pw-sampled-measurable-umf}
Let \((Y,\mu)\) be a standard probability space, let \(K\) be a compact
metrizable group, and let \(D\subseteq Y\) be arbitrary. Then
\[
 M(G_{\mu,D}(K))\cong K^{A_\mu\cup D}.
\]
More precisely, the two-sided completion of \(G_{\mu,D}(K)\) is naturally
\[
 L^0(Y_{\mathrm c},\mu;K)\times K^{A_\mu\cup D},
\]
where an atomic coordinate occurring also in \(D\) is included only once.
\end{theorem}

\begin{proof}
\uses{lem:pw-umf-principles}
An element of \(L^0(Y,\mu;K)\) has a well-defined value on every atom. Hence
the closure of the image of \(\jmath_{\mu,D}\) is the closed subgroup
\[
 \{(g,w)\in L^0(Y,\mu;K)\times K^D:
        w_a=g(a)\text{ for }a\in A_\mu\cap D\}.
\]
It is isomorphic to the displayed product. To prove density, choose a Borel
representative of \(g\). A basic neighbourhood specifies only finitely many
coordinates in \(D\); changing the representative at the non-atomic points
among them does not alter its \(L^0\)-class, while compatibility at the atomic
points is already part of the definition of the closed subgroup.

By \cite[Theorem~1.1]{PestovSchneider2017}, the first factor is extremely amenable. Projection
onto \(K^{A_\mu\cup D}\) therefore has an extremely amenable kernel and compact
quotient. Lemma~\ref{lem:pw-umf-principles}, together with invariance under
two-sided completion, proves the formula.
\end{proof}

For an atomless smooth probability space and connected compact Lie target, the
dense smooth subgroup of \(G_{\mu,D}(K)\) is precisely the group in
Theorem~\ref{thm:pw-mixed-sampling-umf}. Thus
Theorem~\ref{thm:pw-sampled-measurable-umf} both explains that result and adds
arbitrary atoms and arbitrary compact metrizable targets.

We next apply the same extremely-amenable-kernel argument to transformations of
a measure space. Let \((Y,\mu)\) be a standard probability space, retain the
notation \(A_\mu\) for its set of atomic points, and put
\[
Y_{\mathrm c}=Y\setminus A_\mu,
\qquad
R_\mu=\{\mu(\{a\}):a\in A_\mu\},
\qquad
A_r=\{a\in A_\mu:\mu(\{a\})=r\}.
\]
Every \(A_r\) is finite. Write \(\operatorname{Aut}(Y,\mu)_{\mathrm w}\) for the
measure-preserving automorphism group with its weak topology and
\(\operatorname{Aut}^*(Y,\mu)_{\mathrm w}\) for the nonsingular automorphism group with
the corresponding weak topology. Equivalently, the latter is the topology
induced by the quasi-regular representation
\[
 U_\tau^{(p)}f=
 \left(\frac{d(\mu\circ\tau^{-1})}{d\mu}\right)^{1/p}
 f\circ\tau^{-1}
\]
in the strong operator topology of \(L^p(Y,\mu)\), for one (and hence every)
\(1\leq p<\infty\); see \cite[Section~6.2]{GiordanoPestov2007}. This convention will also
be used below for vector-valued quasi-regular representations.

\begin{theorem}\label{thm:pw-measure-aut-umf}
There are natural isomorphisms of flows
\begin{align}
M(\operatorname{Aut}(Y,\mu)_{\mathrm w})
 &\cong \prod_{r\in R_\mu}\operatorname{Sym}(A_r),
 \label{eq:pw-measure-preserving-aut-umf}\\
M(\operatorname{Aut}^*(Y,\mu)_{\mathrm w})
 &\cong \operatorname{LO}(A_\mu),
 \label{eq:pw-nonsingular-aut-umf}
\end{align}
where \(\operatorname{LO}(A_\mu)\) is the compact space of all linear orders
on \(A_\mu\), with the natural permutation action. For a finite atomic set this
flow is the regular translation flow of \(\operatorname{Sym}(A_\mu)\), and for
an empty atomic set it is a point.
\end{theorem}

\begin{proof}
\uses{lem:pw-umf-principles}
A measure-preserving automorphism preserves atom masses, while a nonsingular
automorphism may permute all atoms. Restriction to the atomic and nonatomic
parts gives topological direct-product decompositions
\begin{align*}
\operatorname{Aut}(Y,\mu)_{\mathrm w}
 &\cong \operatorname{Aut}(Y_{\mathrm c},\mu)_{\mathrm w}
   \times\prod_{r\in R_\mu}\operatorname{Sym}(A_r),\\
\operatorname{Aut}^*(Y,\mu)_{\mathrm w}
 &\cong \operatorname{Aut}^*(Y_{\mathrm c},\mu)_{\mathrm w}
   \times\operatorname{Sym}(A_\mu).
\end{align*}
The product topology on the atomic coordinates is the induced weak topology:
control of finitely many positive-mass atoms gives a basic coordinate
condition, and the remaining atomic mass can be made arbitrarily small.

Giordano and Pestov prove that both nonatomic factors above are extremely
amenable, already in \cite[Theorems~2.1 and~2.4]{GiordanoPestov2002}; see also the fuller
proofs in \cite[Theorem~4.2 and Theorem~6.1]{GiordanoPestov2007}. The remaining factor in the
first decomposition is compact, which proves
\eqref{eq:pw-measure-preserving-aut-umf}. If \(A_\mu\) is countably infinite,
Glasner and Weiss compute
\(M(\operatorname{Sym}(A_\mu))=\operatorname{LO}(A_\mu)\)
\cite{GlasnerWeiss2002}. For finite \(A_\mu\), the permutation action on its linear orders
is simply transitive, so it is the compact group's own universal minimal flow.
Lemma~\ref{lem:pw-umf-principles} proves
\eqref{eq:pw-nonsingular-aut-umf} in every case.
\end{proof}


\subsection{Wreath products and atomic decompositions}

Let a nonsingular transformation \(T\) act on \(L^0(Y,\mu;K)\) by
\[
 (T\mathbin{\cdot}f)(y)=f(T^{-1}y).
\]
For a compact metrizable group \(K\), form the measurable wreath products
\begin{align*}
 W_\mu(K)
  &=L^0(Y,\mu;K)\rtimes\operatorname{Aut}(Y,\mu)_{\mathrm w},\\
 W_\mu^*(K)
  &=L^0(Y,\mu;K)\rtimes\operatorname{Aut}^*(Y,\mu)_{\mathrm w}.
\end{align*}

\begin{theorem}\label{thm:pw-measurable-wreath-umf}
There are natural isomorphisms of flows
\begin{align}
 M(W_\mu(K))
 &\cong K^{A_\mu}\rtimes
       \prod_{r\in R_\mu}\operatorname{Sym}(A_r),
 \label{eq:pw-preserving-wreath-umf}\\
 M(W_\mu^*(K))
 &\cong K^{A_\mu}\times\operatorname{LO}(A_\mu).
 \label{eq:pw-nonsingular-wreath-umf}
\end{align}
The first compact group acts on itself by left translation. In the second
formula, the atomic subgroup
\(K^{A_\mu}\rtimes\operatorname{Sym}(A_\mu)\) acts by
\[
 (z,\sigma)\mathbin{\cdot}(w,<)
 =\bigl(z(\sigma\mathbin{\cdot}w),\sigma\mathbin{\cdot}<\bigr),
 \qquad
 (\sigma\mathbin{\cdot}w)_a=w_{\sigma^{-1}(a)},
\]
and the nonatomic factor acts trivially. In particular, both measurable wreath
products are extremely amenable when \(\mu\) is atomless.
\end{theorem}

\begin{proof}
\uses{lem:pw-umf-principles}
Restriction to the atomic and nonatomic parts gives
\begin{align*}
 W_\mu(K)
 &\cong
 \bigl(L^0(Y_{\mathrm c},\mu;K)\rtimes
       \operatorname{Aut}(Y_{\mathrm c},\mu)_{\mathrm w}\bigr)
 \times
 \left(K^{A_\mu}\rtimes
       \prod_{r\in R_\mu}\operatorname{Sym}(A_r)\right),\\
 W_\mu^*(K)
 &\cong
 \bigl(L^0(Y_{\mathrm c},\mu;K)\rtimes
       \operatorname{Aut}^*(Y_{\mathrm c},\mu)_{\mathrm w}\bigr)
 \times
 \bigl(K^{A_\mu}\rtimes\operatorname{Sym}(A_\mu)\bigr).
\end{align*}
Each first factor is extremely amenable: its \(L^0\)-kernel is extremely
amenable by \cite[Theorem~1.1]{PestovSchneider2017}, its quotient is extremely amenable by
\cite[Theorem~4.2 and Theorem~6.1]{GiordanoPestov2007}, and
Lemma~\ref{lem:pw-umf-principles} applies twice.

The second factor in the measure-preserving decomposition is compact, which
gives \eqref{eq:pw-preserving-wreath-umf}. In the nonsingular case,
\(K^{A_\mu}\) is compact and normal. The semidirect-product formula
\cite[Proposition~6.1]{KechrisSokic2012}, together with
\(M(\operatorname{Sym}(A_\mu))=\operatorname{LO}(A_\mu)\) from \cite{GlasnerWeiss2002},
gives the phase space and the stated skew action in
\eqref{eq:pw-nonsingular-wreath-umf}. The extension splits, so no cocycle
correction is needed here: the skew action is the one produced by the
splitting, as in the proof of \cite[Proposition~6.1]{KechrisSokic2012}. The extremely amenable direct factors do
not change either universal minimal flow.
\end{proof}

The compact fibre in the preceding theorem can be replaced by a much larger
class without changing the form of the answer.

\begin{definition}\label{def:pw-ea-by-compact}
A Polish group \(F\) is \emph{extremely-amenable-by-compact} if it has a closed
extremely amenable normal subgroup \(F_0\) such that

\[
 Q_F=F/F_0
\]

is compact.
\end{definition}

\begin{theorem}\label{thm:pw-ea-compact-measurable-wreath}
\uses{thm:pw-measurable-wreath-umf}
Let \(F\) be extremely-amenable-by-compact and give \(L^0(Y,\mu;F)\)
the topology of convergence in measure. Then
\begin{align}
 M(L^0(Y,\mu;F))
  &\cong Q_F^{A_\mu},
  \label{eq:pw-ea-compact-l0-umf}\\
 M\!\left(L^0(Y,\mu;F)\rtimes\operatorname{Aut}(Y,\mu)_{\mathrm w}\right)
  &\cong Q_F^{A_\mu}\rtimes
        \prod_{r\in R_\mu}\operatorname{Sym}(A_r),
  \label{eq:pw-ea-compact-preserving-wreath}\\
 M\!\left(L^0(Y,\mu;F)\rtimes\operatorname{Aut}^*(Y,\mu)_{\mathrm w}\right)
  &\cong Q_F^{A_\mu}\times\operatorname{LO}(A_\mu).
  \label{eq:pw-ea-compact-nonsingular-wreath}
\end{align}
The compact group in \eqref{eq:pw-ea-compact-preserving-wreath} acts on itself by
left translation.  In \eqref{eq:pw-ea-compact-nonsingular-wreath} the quotient
\(Q_F^{A_\mu}\rtimes\operatorname{Sym}(A_\mu)\) has the standard skew action
described in Theorem~\ref{thm:pw-measurable-wreath-umf}, and the remaining
nonatomic factor acts trivially.  In particular, all three groups are
extremely amenable when \(\mu\) is atomless.
\end{theorem}

\begin{proof}
\uses{lem:pw-umf-principles}
The group \(F\) is amenable, since both \(F_0\) and its compact quotient are
amenable.  On the nonatomic part, Pestov and Schneider's theorem therefore
makes \(L^0(Y_{\mathrm c},\mu;F)\) extremely amenable
\cite[Theorem~1.1]{PestovSchneider2017}.  On the atoms, the closed normal subgroup
\(F_0^{A_\mu}\) is extremely amenable and the quotient is the compact group
\(Q_F^{A_\mu}\).  Restriction to the two parts, followed by
Lemma~\ref{lem:pw-umf-principles}, proves \eqref{eq:pw-ea-compact-l0-umf}.

After adjoining measure-space transformations, the nonatomic semidirect
factor is still extremely amenable: its kernel is the preceding measurable-map
group and its quotient is extremely amenable by \cite{GiordanoPestov2007}.  On the atomic
part, divide first by \(F_0^{A_\mu}\).  In the measure-preserving case the
remaining group is the compact group displayed in
\eqref{eq:pw-ea-compact-preserving-wreath}.  In the nonsingular case it is

\[
 Q_F^{A_\mu}\rtimes\operatorname{Sym}(A_\mu).
\]

The compact-normal-subgroup formula \cite[Proposition~6.1]{KechrisSokic2012} and the
Glasner--Weiss computation \cite{GlasnerWeiss2002} give
\eqref{eq:pw-ea-compact-nonsingular-wreath} and its skew action.
\end{proof}

The same compact-extension argument gives a general supply of examples over
countable structures.

\begin{corollary}\label{cor:pw-nonarch-wreath-umf}
Let \(\Omega\) be countable, let
\(H\leq\operatorname{Sym}(\Omega)\) be closed in the topology of pointwise
convergence, and let \(K\) be compact metrizable. Then
\[
 M(K^\Omega\rtimes H)\cong K^\Omega\times M(H),
\]
with the standard skew action. Consequently,
\begin{align*}
 M(K^{\mathbb Q}\rtimes\operatorname{Aut}(\mathbb Q,<))
 &\cong K^{\mathbb Q},\\
 M(K^{\mathbb N}\rtimes\operatorname{Aut}(\mathcal R))
 &\cong K^{\mathbb N}\times\operatorname{LO}(\mathbb N),
\end{align*}
where \(\mathcal R\) is the countable random graph.
\end{corollary}

\begin{proof}
The group \(K^\Omega\) is compact and normal, so the first formula is
\cite[Proposition~6.1]{KechrisSokic2012}. The group
\(\operatorname{Aut}(\mathbb Q,<)\) is extremely amenable, while the universal
minimal flow of \(\operatorname{Aut}(\mathcal R)\) is the space of all linear
orders on its vertex set; both statements follow from the Ramsey-theoretic
computations in \cite{KPT2005}.
\end{proof}

The compactness of the fibre in Corollary~\ref{cor:pw-nonarch-wreath-umf} is
unnecessary when its universal minimal flow is metrizable.  This gives a
wreath-product construction which is not restricted to compact homogeneous
spaces. Before giving its stabiliser proof, we record the earlier comparisons.
Soki\'c's relational composition theorem \cite[Theorem~6.5(v)]{Sokic2013}
identifies the universal flow with admissible convex orders: the fibre class
has strong amalgamation, at least one class has strong joint embedding, and
the reasonable ordered expansions have the Ramsey and ordering properties. With finite
sets as the quotient class, a convex order consists of an order on the
quotient and an admissible order in each fibre; the automorphism group and
flow are respectively \(F^{\mathbb N}\rtimes\operatorname{Sym}(\mathbb N)\) and
\(M(F)^{\mathbb N}\times\operatorname{LO}(\mathbb N)\). Compositions of more general
expansions are treated in \cite[Proposition~2.2 of the 2017 preprint]{PawliukSokic2020}.
Jahel and Zucker's \cite[Theorem~5.1]{JahelZucker2026} instead establishes metrizability
and unique ergodicity for Polish wreath products. The next theorem gives the
explicit flow for arbitrary Polish fibres with metrizable flow, using the
standard co-precompact stabiliser construction; it does not require an ordered
relational presentation of the fibre group.

\begin{theorem}\label{thm:pw-metrizable-fibre-wreath-umf}
Let \(F\) be a Polish group whose universal minimal flow \(X=M(F)\) is
metrizable, and let \(\Omega\) be countable.  Give \(F^\Omega\) the product
topology and \(\operatorname{Sym}(\Omega)\) the pointwise topology.  Then
\[
 M\!\left(F^\Omega\rtimes\operatorname{Sym}(\Omega)\right)
 \cong X^\Omega\times\operatorname{LO}(\Omega).
\]
The action is the standard skew action
\[
 ((f_a)_{a\in\Omega},\sigma)\mathbin{\cdot}((x_a)_{a\in\Omega},<)
 =\bigl((f_a\mathbin{\cdot}x_{\sigma^{-1}(a)})_{a\in\Omega},
        \sigma\mathbin{\cdot}<\bigr).
\]
For finite \(\Omega\), the same formula applies and
\(\operatorname{LO}(\Omega)\) is the simply transitive
\(\operatorname{Sym}(\Omega)\)-space of linear orders.
\end{theorem}

\begin{proof}
By \cite[Theorem~1.1]{BYMT2017}, there are a point \(x_0\in X\) and a closed
extremely amenable co-precompact subgroup
\[
 F^*=\operatorname{Stab}_F(x_0)
\]
such that the completion of \(F/F^*\), for the quotient of the right
uniformity, is \(X\).  If \(\Omega\) is infinite, choose a linear order
\(<_0\) of type \(\mathbb Q\) on \(\Omega\); if it is finite, choose any
linear order.  The subgroup
\[
 H=(F^*)^\Omega\rtimes\operatorname{Aut}(\Omega,<_0)
 \leq F^\Omega\rtimes\operatorname{Sym}(\Omega)
\]
is extremely amenable. Indeed, \((F^*)^\Omega\) is extremely amenable, and
so is the quotient \(\operatorname{Aut}(\Omega,<_0)\), by the Ramsey theorem
in the infinite case and trivially in the finite case. Extreme amenability is
preserved under extensions, so the assertion follows.

Put \(G=F^\Omega\rtimes\operatorname{Sym}(\Omega)\) and
\(y_0=((x_0)_{a\in\Omega},<_0)\). The stabiliser of \(y_0\) is exactly \(H\).
The orbit map
\[
 \Theta:G/H\longrightarrow X^\Omega\times\operatorname{LO}(\Omega),
 \qquad \Theta(gH)=g\mathbin{\cdot}y_0,
\]
is one-to-one, and its image is
\[
 (F\mathbin{\cdot}x_0)^\Omega\times
 \bigl(\operatorname{Sym}(\Omega)\mathbin{\cdot}<_0\bigr).
\]
We spell out the uniform-completion point. A basic entourage of the quotient
right uniformity on \(G/H\) is obtained by choosing a finite set
\(D\subseteq\Omega\), requiring the two induced linear orders to agree on
\(D\),
and requiring, for \(a\in D\), the corresponding two \(F/F^*\)-coordinates
to lie in prescribed entourages of the quotient right uniformity. This follows
from the product topology on \(F^\Omega\), the pointwise topology on
\(\operatorname{Sym}(\Omega)\), and the freedom to multiply on the right by
the independent coordinates of \((F^*)^\Omega\) and by
\(\operatorname{Aut}(\Omega,<_0)\). More explicitly, agreement of the two
induced orders on \(D\) gives an order-preserving bijection between their
inverse images under the two permutations. Homogeneity of the rational order
extends this finite bijection to an order automorphism. After making this
correction, the independent \(F^*\)-coordinates realise each of the finitely
many quotient entourages. The converse follows by restricting a small group
displacement to those coordinates. For finite \(\Omega\), agreement on all
of \(\Omega\) replaces the order-extension step. These are precisely the entourages induced
through \(\Theta\) by the product of the quotient right uniformities on finite
coordinate sets.

By the choice of \(F^*\), the completion of \(F/F^*\) is \(X\). The completion
of \(\operatorname{Sym}(\Omega)/\operatorname{Aut}(\Omega,<_0)\) is
\(\operatorname{LO}(\Omega)\), by the Glasner--Weiss computation (and by
finiteness when \(\Omega\) is finite). The finite-coordinate description
therefore shows that \(H\) is co-precompact and that \(\Theta\) extends to a
\(G\)-isomorphism
\[
 \widehat{G/H}\cong X^\Omega\times\operatorname{LO}(\Omega).
\]
The displayed skew action is minimal: first use a permutation to meet any
prescribed finite order condition and then use the independent \(F\)-coordinates
and minimality of \(X\) to meet any prescribed finite family of open subsets of
\(X\). An extremely amenable co-precompact subgroup makes every minimal
subflow of \(\widehat{G/H}\) a copy of \(M(G)\), as recalled in the introduction
of \cite{BYMT2017}; the minimality just proved therefore gives the formula.
\end{proof}

This produces concrete examples from functional analysis that are not
extremely-amenable-by-compact.  Let \(\mathbb G\) and \(\mathbb{NG}\) denote
the Gurarij Banach space and the noncommutative Gurarij operator space.  Let
\(\mathbb P\) be the Poulsen simplex, and let \(\mathbb{NP}_1\) be the compact
space of scalar states of the noncommutative Poulsen simplex.

\begin{corollary}\label{cor:pw-gurarij-poulsen-wreath-umf}
For every countable \(\Omega\),
\begin{align*}
 M\!\left(\operatorname{Aut}(\mathbb G)^\Omega
          \rtimes\operatorname{Sym}(\Omega)\right)
 &\cong\operatorname{LO}(\Omega),\\
 M\!\left(\operatorname{Aut}(\mathbb{NG})^\Omega
          \rtimes\operatorname{Sym}(\Omega)\right)
 &\cong\operatorname{LO}(\Omega),\\
 M\!\left(\operatorname{Aut}(\mathbb P)^\Omega
          \rtimes\operatorname{Sym}(\Omega)\right)
 &\cong\mathbb P^\Omega\times\operatorname{LO}(\Omega),\\
 M\!\left(\operatorname{Aut}(\mathbb{NP})^\Omega
          \rtimes\operatorname{Sym}(\Omega)\right)
 &\cong\mathbb{NP}_1^\Omega\times\operatorname{LO}(\Omega).
\end{align*}
Here automorphisms mean surjective linear isometries for \(\mathbb G\),
surjective complete isometries for \(\mathbb{NG}\), affine homeomorphisms for
\(\mathbb P\), and matrix-affine homeomorphisms for \(\mathbb{NP}\).
\end{corollary}

\begin{proof}
\uses{thm:pw-metrizable-fibre-wreath-umf}
The groups \(\operatorname{Aut}(\mathbb G)\) and
\(\operatorname{Aut}(\mathbb{NG})\) are extremely amenable
\cite{BartosovaLopezAbadLupiniMbombo2022,BartosovaLopezAbadLupiniMbombo2021}.  The canonical actions on \(\mathbb P\) and
\(\mathbb{NP}_1\) are the respective universal minimal flows of
\(\operatorname{Aut}(\mathbb P)\) and \(\operatorname{Aut}(\mathbb{NP})\),
again by \cite{BartosovaLopezAbadLupiniMbombo2022,BartosovaLopezAbadLupiniMbombo2021}.  Apply
Theorem~\ref{thm:pw-metrizable-fibre-wreath-umf}.
\end{proof}

Theorem~\ref{thm:pw-metrizable-fibre-wreath-umf} and
Corollary~\ref{cor:pw-gurarij-poulsen-wreath-umf} are recorded as consequences of
the cited metrizable-flow and Ramsey results; no separate claim of priority is
made for these formulas.


\subsubsection{Full groups of measured equivalence relations}

Let \(\mathcal R\) be a countable nonsingular measured equivalence relation on
the standard probability space \((Y,\mu)\). Its full group \([\mathcal R]\)
is endowed with the uniform metric
\[
 d_{\mathrm u}(S,T)=\mu\{y\in Y:S(y)\ne T(y)\}.
\]
Put \(\mathcal R_{\mathrm c}=\mathcal R|_{Y_{\mathrm c}}\), and let
\[
 \mathcal C_{\mathcal R}=A_\mu/\mathcal R
\]
be the countable family of \(\mathcal R\)-classes contained in the atomic
part. We regard \(\operatorname{LO}(C)\) as a singleton when \(|C|\leq1\).

\begin{theorem}\label{thm:pw-full-group-atomic-umf}
Assume that \([\mathcal R_{\mathrm c}]\), with the uniform topology, is
extremely amenable. Then
\[
 M([\mathcal R])
 \cong \prod_{C\in\mathcal C_{\mathcal R}}\operatorname{LO}(C),
\]
with the coordinatewise order action. In particular, the conclusion applies
when \(\mathcal R_{\mathrm c}\) is ergodic and amenable. If \(\mathcal R\) is
measure preserving, every atomic class \(C\) is finite, and the displayed flow
is equivariantly isomorphic to the left-translation flow of the compact group
\[
 \prod_{C\in\mathcal C_{\mathcal R}}\operatorname{Sym}(C).
\]
\end{theorem}

\begin{proof}
\uses{lem:pw-umf-principles}
Quasi-invariance prevents an atom from being \(\mathcal R\)-equivalent to a
point of \(Y_{\mathrm c}\): the saturation of a null singleton is null and
therefore cannot contain an atom. Restriction consequently gives a topological
direct-product decomposition
\[
 [\mathcal R]
 \cong [\mathcal R_{\mathrm c}]
 \times\prod_{C\in\mathcal C_{\mathcal R}}\operatorname{Sym}(C).
\]
The topology on the second factor is pointwise convergence. Indeed, fixing a
positive-mass atom is a uniform open condition, while outside a sufficiently
large finite set of atoms the total mass is arbitrarily small.

Write \(H=\prod_C\operatorname{Sym}(C)\). For each countably infinite \(C\),
choose an order \(<_C\) of type \(\mathbb Q\), and choose the unique order type
on each finite \(C\). The subgroup
\[
 H^*=\prod_C\operatorname{Aut}(C,<_C)
\]
is extremely amenable by the Ramsey theorem and closure of extreme amenability
under products. It is co-precompact in \(H\), the completion of \(H/H^*\) is
\(\prod_C\operatorname{LO}(C)\), and the resulting action is minimal. The
co-precompact-subgroup description of universal minimal flows
\cite[Theorem~1.1]{BYMT2017} therefore gives
\[
 M(H)=\prod_C\operatorname{LO}(C);
\]
equivalently, this is the coordinatewise product of the Glasner--Weiss
calculation \cite{GlasnerWeiss2002}. The extremely amenable first direct factor does not
change the universal minimal flow by Lemma~\ref{lem:pw-umf-principles}.

If \(\mathcal R\) is measure preserving, atoms in the same class have equal
mass. Such a class is finite, and the action of \(\operatorname{Sym}(C)\) on
\(\operatorname{LO}(C)\) is simply transitive. Finally, Giordano and Pestov
prove the asserted extreme amenability when \(\mathcal R_{\mathrm c}\) is
ergodic and amenable \cite[Theorem~5.7]{GiordanoPestov2007}.
\end{proof}

The same calculation allows compact labels along the equivalence relation.

\begin{corollary}\label{cor:pw-decorated-full-group-umf}
\uses{thm:pw-full-group-atomic-umf}
Let \(K\) be compact metrizable and retain the assumption of
Theorem~\ref{thm:pw-full-group-atomic-umf}. For the action by precomposition, put
\[
 W_{\mathcal R}(K)=L^0(Y,\mu;K)\rtimes[\mathcal R].
\]
Then
\[
 M(W_{\mathcal R}(K))
 \cong K^{A_\mu}\times
       \prod_{C\in\mathcal C_{\mathcal R}}\operatorname{LO}(C).
\]
The atomic group
\[
 K^{A_\mu}\rtimes
 \prod_{C\in\mathcal C_{\mathcal R}}\operatorname{Sym}(C)
\]
acts by the standard skew action, and the nonatomic factor acts trivially.
\end{corollary}

\begin{proof}
\uses{thm:pw-full-group-atomic-umf}
Restriction to \(Y_{\mathrm c}\) and \(A_\mu\) decomposes
\(W_{\mathcal R}(K)\) as the direct product of
\[
 L^0(Y_{\mathrm c},\mu;K)\rtimes[\mathcal R_{\mathrm c}]
\]
and the displayed atomic group. The first factor is extremely amenable: its
measurable-map kernel is extremely amenable by \cite[Theorem~1.1]{PestovSchneider2017}, its
quotient is extremely amenable by assumption, and fixed points lift through
extensions. On the atomic factor, \(K^{A_\mu}\) is compact and normal.
Theorem~\ref{thm:pw-full-group-atomic-umf} and the compact semidirect-product
formula \cite[Proposition~6.1]{KechrisSokic2012} give the result.
\end{proof}

\begin{corollary}\label{cor:pw-ea-compact-full-group-umf}
\uses{thm:pw-full-group-atomic-umf}
Let \(F\) be extremely-amenable-by-compact and retain the hypothesis on
\(\mathcal R\) from Theorem~\ref{thm:pw-full-group-atomic-umf}. Then
\[
 M\!\left(L^0(Y,\mu;F)\rtimes[\mathcal R]\right)
 \cong Q_F^{A_\mu}\times
       \prod_{C\in\mathcal C_{\mathcal R}}\operatorname{LO}(C).
\]
The quotient
\[
 Q_F^{A_\mu}\rtimes
 \prod_{C\in\mathcal C_{\mathcal R}}\operatorname{Sym}(C)
\]
acts by the coordinatewise standard skew action, and the extremely amenable
kernel acts trivially.
\end{corollary}

\begin{proof}
\uses{thm:pw-ea-compact-measurable-wreath, thm:pw-full-group-atomic-umf}
The nonatomic semidirect factor is extremely amenable, as in the proof of
Theorem~\ref{thm:pw-ea-compact-measurable-wreath}.  On the atomic part, divide by
the extremely amenable normal subgroup \(F_0^{A_\mu}\), then apply the compact
semidirect-product formula \cite[Proposition~6.1]{KechrisSokic2012} and
Theorem~\ref{thm:pw-full-group-atomic-umf}.
\end{proof}

The nonatomic extreme-amenability input in
Theorem~\ref{thm:pw-full-group-atomic-umf} and
Corollaries~\ref{cor:pw-decorated-full-group-umf} and
\ref{cor:pw-ea-compact-full-group-umf} is due to Giordano and Pestov. Retaining
the atoms exposes the order flows of the individual atomic equivalence
classes.

\medskip
For scalar \(\ell^p\), the following atomic skew flow is already stated
in \cite[Theorem~4.13 and Remark~4.14]{PestovRemarks2002}. Splitting the atomic and
nonatomic bands gives the corresponding mixed-space formula. We record the
calculation for the isometry groups of \(L^p\)-spaces. Put
\[
\mathbb S_{\mathbb C}=\mathbb T,
\qquad
\mathbb S_{\mathbb R}=\{-1,1\}.
\]

\begin{theorem}\label{thm:pw-lp-isometry-umf}
Let \(\mathbb F\in\{\mathbb R,\mathbb C\}\) and
\(1\leq p<\infty\), \(p\ne2\). Equip the group of surjective linear
isometries of \(L^p(Y,\mu;\mathbb F)\) with the strong operator topology. Then
the phase space of its universal minimal flow is
\[
\mathbb S_{\mathbb F}^{A_\mu}\times\operatorname{LO}(A_\mu).
\]
The action is the standard skew action of the atomic monomial group
\[
\mathbb S_{\mathbb F}^{A_\mu}\rtimes
\operatorname{Sym}(A_\mu),
\]
acting as follows: if \(z,w\in\mathbb S_{\mathbb F}^{A_\mu}\),
\(\sigma\in\operatorname{Sym}(A_\mu)\), and
\(<\in\operatorname{LO}(A_\mu)\), then
\[
 (z,\sigma)\mathbin{\cdot}(w,<)
 =\bigl(z(\sigma\mathbin{\cdot}w),\sigma\mathbin{\cdot}<\bigr),
 \qquad
 (\sigma\mathbin{\cdot}w)_a=w_{\sigma^{-1}(a)},
\]
where \(a\,(\sigma\mathbin{\cdot}<)\,b\) means
\(\sigma^{-1}(a)<\sigma^{-1}(b)\). The isometries of the nonatomic summand
act trivially on this flow. In
particular, the isometry group is extremely amenable when \(\mu\) is atomless.
\end{theorem}

\begin{proof}
\uses{lem:pw-umf-principles, thm:pw-measure-aut-umf}
The Banach--Lamperti description, in the topological form used in
\cite[Section~6.2]{GiordanoPestov2007}, preserves the atomic and nonatomic bands and gives
\[
\operatorname{Iso}(L^p(Y,\mu;\mathbb F))
\cong E_{\mathrm c}\times
 \bigl(\mathbb S_{\mathbb F}^{A_\mu}\rtimes
       \operatorname{Sym}(A_\mu)\bigr).
\]
Here \(E_{\mathrm c}\) is the isometry group of the nonatomic summand and is
extremely amenable by \cite[Theorem~6.6]{GiordanoPestov2007}. On the atomic summand, after
normalising the characteristic function of every atom, the strong operator
topology is the product topology on the scalar coordinates together with
pointwise convergence on the permutation coordinate.

The group \(\mathbb S_{\mathbb F}^{A_\mu}\) is compact and normal. The
semidirect-product formula for universal minimal flows with a compact normal
factor \cite[Proposition~6.1]{KechrisSokic2012}, together with Theorem~\ref{thm:pw-measure-aut-umf},
identifies the phase space with
\(\mathbb S_{\mathbb F}^{A_\mu}\times\operatorname{LO}(A_\mu)\) and supplies
the stated skew action. The extremely amenable direct factor does not change
the universal minimal flow by Lemma~\ref{lem:pw-umf-principles}.
\end{proof}


\subsection{Operator-algebraic applications}


\subsubsection{Weak unitary groups of injective and nuclear algebras}

The operator-algebraic compact factor can be identified explicitly.
The underlying compact-times-extremely-amenable decomposition is already
\cite[Corollary~3.6]{GiordanoPestov2002}; we record its factor coordinates using the fuller
proof in \cite{GiordanoPestov2007}. Let \(N\)
be an injective von Neumann algebra and let
\(\{z_i:i\in I_{\mathrm{fd}}(N)\}\) be the family of minimal central
projections for which \(z_iN\) is finite-dimensional. Write
\(z_iN\cong M_{n_i}(\mathbb C)\) and put
\[
\mathcal C_N=\prod_{i\in I_{\mathrm{fd}}(N)}\mathrm U(n_i).
\]

\begin{theorem}\label{thm:pw-injective-vn-umf}
For the \(s(N,N_{\ast})\)-topology,
\[
M(\mathrm U(N)_{\mathrm s})\cong\mathcal C_N.
\]
Thus \(\mathrm U(N)_{\mathrm s}\) is extremely amenable precisely when \(N\) has no
nonzero finite-dimensional factor summand.
\end{theorem}

\begin{proof}
\uses{lem:pw-umf-principles}
Let \(z=\sum_{i\in I_{\mathrm{fd}}(N)}z_i\). The central decomposition
\[
N=zN\oplus(1-z)N
\]
gives a topological direct product decomposition of the strong-topological
unitary group. The first factor is the compact group \(\mathcal C_N\).
Giordano and Pestov's proof of \cite[Theorem~3.3]{GiordanoPestov2007} shows that the unitary
group of the complementary summand is extremely amenable: the nonatomic
finite type-I, finite continuous, and properly infinite injective parts are
handled in their Proposition~3.4, Proposition~3.7, and Corollary~3.6,
respectively. Lemma~\ref{lem:pw-umf-principles} gives the formula.
\end{proof}

For example, if
\[
N=L^\infty(Y,\mu)\,\overline\otimes\,M_n(\mathbb C),
\]
then \(\mathrm U(N)_{\mathrm s}=L^0(Y,\mu;\mathrm U(n))\)
\cite[Lemma~3.2]{GiordanoPestov2007}, and Theorem~\ref{thm:pw-injective-vn-umf} recovers
\[
M(\mathrm U(N)_{\mathrm s})\cong\mathrm U(n)^{A_\mu}.
\]

The preceding theorem also gives a formula directly at the \(C^*\)-algebra
level. For a unital \(C^*\)-algebra \(A\), let
\(\widehat A_{\mathrm{fd}}\) be the set of unitary-equivalence classes of
finite-dimensional irreducible representations, and choose a representative
\(\pi:A\to\mathcal B(H_\pi)\) in each class.

\begin{lemma}\label{lem:pw-finite-dimensional-envelope-summands}
Let \(A\) be a unital \(C^*\)-algebra. The finite-dimensional factor summands
of its universal von Neumann envelope \(A^{**}\) are in canonical bijection
with \(\widehat A_{\mathrm{fd}}\). If \(z_{\mathrm{fd}}\) is their central
sum, then
\[
 z_{\mathrm{fd}}A^{**}
 \cong
 \bigoplus_{[\pi]\in\widehat A_{\mathrm{fd}}}^{\ell^\infty}
 \mathcal B(H_\pi),
\qquad
 \mathrm U(z_{\mathrm{fd}}A^{**})_{\mathrm s}
 \cong
 \prod_{[\pi]\in\widehat A_{\mathrm{fd}}}\mathrm U(H_\pi),
\]
where the last group has the product topology.
\end{lemma}

\begin{proof}
A finite-dimensional irreducible representation \(\pi\) has a normal
surjective extension
\(\widetilde\pi:A^{**}\to\mathcal B(H_\pi)\). Its kernel is a weak-star
closed ideal, hence has the form \((1-z_\pi)A^{**}\) for a central projection
\(z_\pi\), and
\(z_\pi A^{**}\cong\mathcal B(H_\pi)\). In particular, \(z_\pi\) is a
minimal nonzero central projection. Hence two such central projections are
either orthogonal or equal. If two such central summands agree,
the corresponding finite-dimensional quotient representations differ by an
automorphism of a full matrix algebra and are therefore unitarily equivalent.
Thus inequivalent irreducible representations give orthogonal summands.

Conversely, restricting a finite-dimensional factor-summand representation of
\(A^{**}\) to the weak-star dense copy of \(A\) gives an irreducible
finite-dimensional representation: in finite dimension weak-star density is
norm density, so the image is the full matrix algebra. These two constructions
are inverse. The first displayed decomposition follows by taking the central
sum. The strong-topological unitary group of a von Neumann direct sum is the
product of the strong-topological unitary groups of its summands
\cite[Lemma~3.1]{GiordanoPestov2007}; on each finite-dimensional summand this is the usual
compact topology.
\end{proof}

\begin{theorem}\label{thm:pw-nuclear-weak-umf}
Let \(A\) be a unital separable nuclear \(C^*\)-algebra. Equip
\(\mathrm U(A)\) with the relative weak topology \(\sigma(A,A^*)\). Then
\[
M(\mathrm U(A)_{\mathrm w})
 \cong
 \prod_{[\pi]\in\widehat A_{\mathrm{fd}}}\mathrm U(H_\pi),
\]
where \(u\in\mathrm U(A)\) acts through the coordinates \(\pi(u)\).
In particular, \(\mathrm U(A)_{\mathrm w}\) is extremely amenable when
\(A\) has no finite-dimensional irreducible representations.
\end{theorem}

\begin{proof}
\uses{lem:pw-finite-dimensional-envelope-summands, thm:pw-injective-vn-umf}
The two-sided completion of \(\mathrm U(A)\) in the relative weak topology is
\(\mathrm U(A^{**})\) with the \(s(A^{**},A^*)\)-topology; see the proof of
\cite[Theorem~3.8]{GiordanoPestov2007}. Nuclearity makes \(A^{**}\) injective. Its
finite-dimensional factor summands are indexed precisely by
\(\widehat A_{\mathrm{fd}}\), by
Lemma~\ref{lem:pw-finite-dimensional-envelope-summands}, and their unitary groups
form the displayed compact product. Theorem~\ref{thm:pw-injective-vn-umf} computes the universal
minimal flow of the completion. Universal minimal flows are unchanged on
passing to the two-sided completion by \cite[Fact~2.3]{BassoZucker2021}.
\end{proof}

\begin{corollary}\label{cor:pw-nuclear-weak-examples}
\uses{prop:pw-nccw-umf}
Let \(X\) be compact metrizable.
\begin{enumerate}
\item
\(M(\mathrm U(C(X))_{\mathrm w})\cong\mathbb T^X\).
\item If \(A\) is unital, separable, nuclear, and has no finite-dimensional
representations, then \(M(\mathrm U(A)_{\mathrm w})\) is a point.
\item For the one-dimensional NCCW algebra in
Proposition~\ref{prop:pw-nccw-umf},
\[
M(\mathrm U(A)_{\mathrm w})\cong\mathcal B_A.
\]
Thus \(\mathcal B_A\) is the full universal minimal flow in both the
pointwise and relative weak topologies, but remains only a known factor in the
norm topology.
\end{enumerate}
\end{corollary}

\begin{proof}
\uses{thm:pw-nuclear-weak-umf}
The irreducible representations of \(C(X)\) are its characters, indexed by
\(X\). The second assertion is the empty-product case of
Theorem~\ref{thm:pw-nuclear-weak-umf}. For a one-dimensional NCCW algebra, the
finite-dimensional irreducible representations consist of the interior fibre
representations together with the representations of the finite-dimensional
quotient \(E\). Consequently, the product in
Theorem~\ref{thm:pw-nuclear-weak-umf} is naturally
\(\mathrm U(E)\times\mathrm U(F)^{(0,1)}\), which is isomorphic to
\(\mathcal B_A\) by filling in the endpoint coordinates using
\(\phi_{\mathrm U}\).
\end{proof}

There is a corresponding formula after dividing by scalar unitaries. Put
\[
 \operatorname{PU}(A)=\mathrm U(A)/(\mathbb T1_A),
 \qquad
 \mathcal Q_A=
 \prod_{[\pi]\in\widehat A_{\mathrm{fd}}}\mathrm U(H_\pi),
\]
and, when the product is nonempty, let
\[
 \Delta\mathbb T
 =\{(\lambda1_{H_\pi})_{[\pi]}:\lambda\in\mathbb T\}
 \leq\mathcal Q_A.
\]
For the empty product, \(\mathcal Q_A/\Delta\mathbb T\) is understood to be a
point.

We use the following elementary completion fact. If \(G\) is a Hausdorff
topological group with two-sided completion \(\widehat G\), and \(C\leq G\)
is compact and central, then the two-sided completion of \(G/C\) is
\(\widehat G/C\). Indeed, compactness makes \(C\) closed in \(\widehat G\);
the quotient uniformity on \(G/C\) is the restriction of that on
\(\widehat G/C\), and the image of \(G/C\) is dense. This also shows that no
additional kernel appears when the scalar circle is removed below.

\begin{theorem}\label{thm:pw-projective-nuclear-weak-umf}
Let \(A\) be a unital separable nuclear \(C^*\)-algebra, and equip
\(\operatorname{PU}(A)\) with the quotient of the relative weak topology on
\(\mathrm U(A)\). Then
\[
 M(\operatorname{PU}(A)_{\mathrm w})
 \cong \mathcal Q_A/\Delta\mathbb T.
\]
The action is induced coordinatewise by the finite-dimensional irreducible
representations: the class of \(u\in\mathrm U(A)\) sends the class of
\((v_\pi)_{[\pi]}\) to the class of
\((\pi(u)v_\pi)_{[\pi]}\).
\end{theorem}

\begin{proof}
\uses{lem:pw-umf-principles, thm:pw-injective-vn-umf, thm:pw-nuclear-weak-umf}
Let \(z\in A^{**}\) be the sum of the finite-dimensional factor summands and
set
\[
 E_A=\mathrm U((1-z)A^{**}).
\]
The proof of Theorem~\ref{thm:pw-nuclear-weak-umf} identifies the two-sided
completion of \(\mathrm U(A)_{\mathrm w}\) with
\(E_A\times\mathcal Q_A\), and Theorem~\ref{thm:pw-injective-vn-umf} shows that
\(E_A\) is extremely amenable. Since the scalar circle is compact and central,
the completion of \(\operatorname{PU}(A)_{\mathrm w}\) is
\[
 (E_A\times\mathcal Q_A)/\Delta_A\mathbb T,
 \qquad
 \Delta_A\mathbb T
 =\{(\lambda1_{(1-z)A^{**}},
       (\lambda1_{H_\pi})_{[\pi]}):\lambda\in\mathbb T\}.
\]
If \(\mathcal Q_A\) is nontrivial, projection onto the second coordinate
induces a quotient onto \(\mathcal Q_A/\Delta\mathbb T\) whose kernel is
isomorphic to \(E_A\). If \(\mathcal Q_A\) is trivial, the completion is
\(E_A/(\mathbb T1)\), which is extremely amenable because quotients of
extremely amenable groups are extremely amenable. Thus in either case the
displayed compact quotient has an extremely amenable closed normal kernel.
Lemma~\ref{lem:pw-umf-principles}, followed by invariance of the universal minimal
flow under two-sided completion \cite[Fact~2.3]{BassoZucker2021}, proves the formula.
\end{proof}

\begin{corollary}\label{cor:pw-projective-nuclear-weak-examples}
\uses{prop:pw-nccw-umf}
Let \(X\) be compact metrizable.
\begin{enumerate}
\item
\[
 M(\operatorname{PU}(C(X))_{\mathrm w})
 \cong\mathbb T^X/\Delta\mathbb T.
\]
\item If \(A\) is unital, separable, nuclear, and has no finite-dimensional
irreducible representations, then \(\operatorname{PU}(A)_{\mathrm w}\) is
extremely amenable.
\item For the one-dimensional NCCW algebra in
Proposition~\ref{prop:pw-nccw-umf},
\[
 M(\operatorname{PU}(A)_{\mathrm w})
 \cong\mathcal B_A/\Delta\mathbb T.
\]
\end{enumerate}
\end{corollary}

\noindent\textit{Proof.}
For \(C(X)\), the finite-dimensional irreducible representations are the
characters indexed by \(X\). The second assertion is the empty-product case of
Theorem~\ref{thm:pw-projective-nuclear-weak-umf}. The identification
\(\mathcal Q_A\cong\mathcal B_A\) in the proof of
Corollary~\ref{cor:pw-nuclear-weak-examples} intertwines the two diagonal scalar
circles, which gives the last formula.
\hfill\(\square\)
\par

The weak-topological projective unitary formula is a consequence of the
preceding completion argument. For \(A=M_n(\mathbb C)\),
the formula reduces, as it should, to the compact group
\(\operatorname{PU}(n)=\mathrm U(n)/\mathbb T\).

\subsection{Proofs of the statements in Section~\ref{sec:maps-transformations}}\label{sec:pw-encyclopaedia-proofs}

\begin{proof}[Proof of Theorem~\ref{thm:lupini-cp}]
\proves{thm:lupini-cp}
\uses{thm:pw-cp-umf}
This is the first assertion of Theorem~\ref{thm:pw-cp-umf}.
\end{proof}

\begin{proof}[Proof of Theorem~\ref{thm:lupini-aut-measure}]
\proves{thm:lupini-aut-measure}
\uses{thm:pw-measure-aut-umf}
This is Theorem~\ref{thm:pw-measure-aut-umf}.
\end{proof}

\begin{proof}[Proof of Theorem~\ref{thm:lupini-lp}]
\proves{thm:lupini-lp}
\uses{thm:pw-lp-isometry-umf}
This is Theorem~\ref{thm:pw-lp-isometry-umf}.
\end{proof}

\begin{proof}[Proof of Theorem~\ref{thm:lupini-wreath}]
\proves{thm:lupini-wreath}
\uses{thm:pw-metrizable-fibre-wreath-umf}
This is Theorem~\ref{thm:pw-metrizable-fibre-wreath-umf}.
\end{proof}

\begin{proof}[Proof of Theorem~\ref{thm:lupini-nuclear}]
\proves{thm:lupini-nuclear}
\uses{thm:pw-nuclear-weak-umf, cor:pw-nuclear-weak-examples}
The formula and the extreme amenability are Theorem~\ref{thm:pw-nuclear-weak-umf};
the formula for $U(C(X))$ is part~(1) of Corollary~\ref{cor:pw-nuclear-weak-examples}.
\end{proof}
